\chapter{Your First Magical Character}
\label{ch:first-magical}
\index{magic!first character}

So you want to play a magic‑user. Good choice – but \textsc{Fate’s Edge} offers \textbf{six different magical paths}. Jumping into all of them at once is a recipe for confusion. This guide walks you through building a \textbf{single, focused magical character} that is powerful, fun, and manageable.

\section{Step 0: The Recommendation}
\index{Runekeeper!as first character}

If this is your first magical character, play a \textbf{Runekeeper}. Why?

\begin{itemize}
  \item \textbf{One Patron:} You swear to a single Patron. No juggling multiple Symbols or conflicting obligations.
  \item \textbf{Clear costs:} Every Rite costs Obligation. Track one number.
  \item \textbf{Support from your Thiasos:} Your familiar is a built‑in roleplay hook and mechanical aid.
\end{itemize}

You can explore Invokers, Cantors, Witches, Psions, or Summoners later. Start simple.

\section{Step 1: Choose a Patron}
\index{Patrons!choosing a first Patron}

Your Patron defines your Rites, your Obligation, and your role in the story. Pick one that matches your character concept.

\begin{center}
\small
\renewcommand{\arraystretch}{1.2}
\begin{tabular}{lll}
\toprule
\textbf{Patron} & \textbf{Domain} & \textbf{Character Concept} \\
\midrule
The Traveler & Roads, journeys, safe passage & Guide, courier, scout \\
Oath of Flame \& Light & Dawn, vows, protection & Paladin, healer, knight \\
Grimmir & Nature, seasons & Druid, ranger, hedge‑witch \\
Ikasha & Shadow, secrets, thresholds & Spy, infiltrator, night‑walker \\
The Witness & Truth, memory, revelation & Investigator, judge, archivist \\
Mykkiel & Law, covenant, judgment & Templar, inquisitor, notary \\
The Carrion‑King & Decay, renewal, compost & Necromancer (gentle), healer \\
\bottomrule
\end{tabular}
\end{center}

\textbf{New player favourite:} \emph{The Traveler}. Low Obligation costs, useful Rites, and a forgiving Patron.

\section{Step 2: Buy Your Core Talents}
\index{talents!Runekeeper requirements}

At character creation (30–34 XP), you need two talents:

\begin{itemize}
  \item \textbf{Familiar (Thiasos) – 2 XP.} Your bond to your Patron. It might be a raven (Traveler), a candle flame (Oath of Flame), or a small fungal creature (Carrion‑King).
  \item \textbf{Codex – 4 XP.} Your book of Rites. It could be a leather journal, a set of carved slate tablets, or a roll of bark.
\end{itemize}

\textbf{Total so far: 6 XP.} You have 24–28 XP remaining for Attributes, Skills, and optional talents.

\section{Step 3: Set Your Attributes}
\index{Attributes!for Runekeepers}

Runekeepers rely on \textbf{Spirit} (Obligation capacity, Rite effectiveness) and \textbf{Wits} (Lore, Investigation). A simple spread:

\begin{center}
\begin{tabular}{lll}
\toprule
\textbf{Attribute} & \textbf{Rating} & \textbf{Cost (new rating ×3)} \\
\midrule
Spirit & 3 & 9 XP (1→2 = 6, 2→3 = 9) \\
Wits & 2 & 6 XP (1→2 = 6) \\
Body & 1 & 0 (base) \\
Presence & 1 & 0 (base) \\
\bottomrule
\end{tabular}
\end{center}

\textbf{Total: 15 XP.} You have 9–13 XP remaining.

\section{Step 4: Buy Your Skills}
\index{Skills!for Runekeepers}

You need \textbf{Lore} (to understand your Rites) and at least one social or utility skill.

\begin{center}
\begin{tabular}{lll}
\toprule
\textbf{Skill} & \textbf{Rating} & \textbf{Cost (new level ×2)} \\
\midrule
Lore & 2 & 6 XP (0→1 = 2, 1→2 = 4) \\
Insight & 1 & 2 XP (0→1 = 2) \\
Sway & 1 & 2 XP (0→1 = 2) \\
\bottomrule
\end{tabular}
\end{center}

\textbf{Total: 10 XP.} You now have 0–3 XP left – perfect for one minor talent.

\section{Step 5: Choose Your First Rites}
\index{Rites!starting Rites}

You begin with \textbf{two Low Rites} from your Patron. Pick ones that are useful and clear.

\subsection*{Example: The Traveler}
\begin{itemize}
  \item \textbf{Road‑Sense (Low, 4 XP):} Gain +1 die to avoid ambushes or delays on a journey. \textit{Use this every travel leg.}
  \item \textbf{Traveler’s Boon (Low, 5 XP):} Ignore one level of difficult terrain or bureaucracy for the scene. \textit{Great for crossing rivers or skipping paperwork.}
\end{itemize}

\subsection*{Example: Oath of Flame \& Light}
\begin{itemize}
  \item \textbf{Kindle Vow (Low, 4 XP):} Declare a short vow (e.g., “I will protect the villagers”). Gain +1 die to actions fulfilling it.
  \item \textbf{Lay on Hands (Low, 5 XP):} Cleanse an affliction or downgrade Harm by 1. \textit{The classic healing Rite.}
\end{itemize}

\section{Step 6: Understand Your Resources}
\index{Obligation!tracking}
\index{Fatigue!tracking}

You track three numbers:

\begin{enumerate}
  \item \textbf{Obligation (to your Patron).} Capacity = Spirit + Presence. For Spirit 3, Presence 1, capacity is 4. Each Rite marks +1 Obligation. When you exceed capacity, mark 1 Fatigue per extra segment.
  \item \textbf{Fatigue.} Each point worsens Position. When Fatigue fills your Body track (Body 1 = 1 segment), convert to Harm.
  \item \textbf{Harm.} Harm 1 = –1 die, Harm 2 = –2 dice, Harm 3 = incapacitated.
\end{enumerate}

\textbf{Playerʼs job:} Mark these yourself. Announce when Obligation fills. The GM does not track them for you.

\section{Step 7: Play Your First Session}
\label{sec:first-session-magic}

\subsection*{Before the session}
\begin{itemize}
  \item Write your two Rites on an index card.
  \item Set your Obligation to 0.
  \item Set your Fatigue to 0.
\end{itemize}

\subsection*{During the session}
\begin{enumerate}
  \item \textbf{Invoke a Rite:} Say “I invoke \emph{Road‑Sense}.” Mark +1 Obligation. Roll Lore + Spirit (or as specified). Apply the effect.
  \item \textbf{Take a Short Rest:} Clear 2 Fatigue.
  \item \textbf{Accept a Patron’s demand:} If the GM says “your Patron requires a task,” treat it as a story hook, not a punishment.
\end{enumerate}

\subsection*{After the session}
\begin{itemize}
  \item Clear Obligation by performing an \textbf{Act of Service} for your Patron (e.g., The Traveler: guide a lost traveller; Oath of Flame: speak a public vow). Describe the act. The GM will tell you how much Obligation clears.
  \item Spend XP to raise Attributes, Skills, or buy new Rites.
\end{itemize}

\section{Optional: Add a Minor Talent}
\index{talents!minor}

If you have 2–4 XP left, take one of these:

\begin{itemize}
  \item \textbf{Efficient Invocation (4 XP):} Reduce Obligation cost of one specific Low Rite by 1 (minimum 0).
  \item \textbf{Keen Senses (2 XP):} +1 die on perception checks.
  \item \textbf{Silver Tongue (2 XP):} +1 die on persuasion attempts.
\end{itemize}

\section{What to Avoid at First}
\label{sec:avoid-first}

\begin{itemize}
  \item \textbf{Multiple Patrons.} Stick to one. Cross‑Resonance (page XX) adds complexity you don’t need yet.
  \item \textbf{Crack the Seal (Invoker ability).} You are a Runekeeper. Let Invokers handle emergency rituals.
  \item \textbf{Cantor’s Corruption.} A great path, but the Corruption clock requires active management. Save it for your second character.
\end{itemize}

\section{Sample Build: Kestra the Traveler’s Runekeeper}
\label{sec:sample-build}

\begin{center}
\small
\renewcommand{\arraystretch}{1.1}
\begin{tabular}{ll}
\toprule
\textbf{Element} & \textbf{Choice} \\
\midrule
Name & Kestra \\
Patron & The Traveler \\
Attributes & Spirit 3 (9 XP), Wits 2 (6 XP), Body 1, Presence 1 \\
Skills & Lore 2 (6 XP), Insight 1 (2 XP), Sway 1 (2 XP) \\
Talents & Familiar (2 XP), Codex (4 XP) \\
Rites & Road‑Sense (4 XP), Traveler’s Boon (5 XP) \\
Starting Boons & 3 \\
Obligation Capacity & 4 (Spirit 3 + Presence 1) \\
\bottomrule
\end{tabular}
\end{center}

\textbf{Total XP spent:} 9+6+2+2+2+4+4+5 = 34 XP (using the +4 XP from Bonds/Complications).

\medskip
\noindent Kestra is a former caravan guide who swore to the Traveler after being betrayed by her crew. She carries a brass compass (her Thiasos) and a leather journal of route‑maps (her Codex). She invokes \emph{Road‑Sense} before every journey and \emph{Traveler’s Boon} to cross difficult terrain. She tracks her Obligation carefully and pays it off by guiding lost travellers through the mountains.

\section{Final Advice}
\label{sec:final-advice-magic}

\begin{tcolorbox}[
  colback=green!5!white,
  colframe=green!75!black,
  title=The One Rule for New Magical Players,
  fonttitle=\bfseries
]
\emph{Spend your Boons. Hoarding them for XP conversion starves the table of drama. A re‑rolled failure, an improved Position, or an activated Asset creates a memorable moment. The XP will come from your adventures.}
\end{tcolorbox}

You are now ready to play a Runekeeper. When you feel comfortable, read about the other five paths (Chapter 3). But for now, enjoy the clarity of one Patron, one Codex, and one set of Rites.

The road is waiting. Pay your toll.